\section{Commitment-Layer Composition}
\label{sec:commitment}

\subsection{Motivation}

The revealed-sacrifice channel observes trade events.  In its raw
form, each event exposes $(i, X, Y, t)$---the agent's identity,
what they surrendered, what they acquired, and when.  This is
already more private than direct $\volR$ measurement, but the
framework can do better: the commitment protocol
of~\citep[Definition~6]{lasser2026exo} allows the sacrifice event
to be committed with hiding, so that only the $\volR$-relevant
content (the magnitude $\Delta\volL(X)$ and the category of $Y$)
is revealed.

\subsection{Formal construction}

\begin{definition}[Committed Revealed-Sacrifice Event]
\label{def:committed_event}
A committed revealed-sacrifice event is a tuple
$(\Gamma, \pi, t)$ where:
\begin{itemize}
\item $\Gamma = \Commit((\Delta\volL(X), \catY), r)$ is a
  commitment to the sacrifice magnitude and bundle category,
  using randomness~$r$~\citep[Definition~6]{lasser2026exo};
\item $\pi$ is a zero-knowledge
  proof~\citep[Definition~7]{lasser2026exo}
  that:
  \begin{enumerate}
  \item[(a)] the commitment $\Gamma$ corresponds to a valid
    revealed-sacrifice event satisfying the objectively
    verifiable components: S3 (bundle coherence) and the
    objective component of S2;
  \item[(b)] $\Delta\volL(X) \geq 0$;
  \item[(c)] the trade actually occurred (linked to a verifiable
    ledger entry);
  \item[(d)] the bundle $Y$ lies wholly within the claimed
    partition component: $Y \subseteq C_{\catY}$
    (Definition~\ref{def:category_partition}(ii));
  \end{enumerate}
\item $t$ is the event timestamp (public).
\end{itemize}

The committed event reveals $\Delta\volL(X)$ (the sacrifice
magnitude) and $\catY$ (the bundle's partition label).  It hides
agent identity~$i$, the specific benchmarked capability~$X$, the
specific unbenchmarked bundle~$Y$, and the counterparty.
\end{definition}

\begin{remark}[S0, S1, S4(a), and S5 are not ZK-verifiable]
\label{rem:zk_limits}
Assumptions S0 (calibration), S1 (free choice), S4(a) (valuation
additivity), and S5 (decomposition validity) are about the
agent's internal state or the true $\volR$ distribution.  These
cannot be verified in zero
knowledge~\citep{goldwasser1985knowledge} without an oracle for
the agent's subjective state.  The ZK proof verifies the
\emph{objective} conditions (S2, S3, and S4(b) via poset
structure); the \emph{subjective} conditions are structural
assumptions about the trade environment, not properties the
commitment protocol can enforce.
\end{remark}

\begin{definition}[Category Partition]
\label{def:category_partition}
The function $\catY$ maps each unbenchmarked bundle to a label in
a \emph{public partition}
$\mathcal{C} = \{C_1, \ldots, C_m\}$ of the unbenchmarked
capability space $U$, subject to two conditions:
\begin{enumerate}
\item[(i)] \emph{Poset-independence:} for any $j \neq k$, no
  capability in $C_j$ is related by subsumption or cooperative
  composition to any capability in $C_k$.
\item[(ii)] \emph{Bundle containment:} each bundle $Y$ lies
  wholly within a single partition component:
  $\catY = j$ iff $Y \subseteq C_j$.
\end{enumerate}
The partition is published as part of the framework's measurement
infrastructure and is fixed for a given trade window.  This
definition is load-bearing for aggregation: distinct category
labels imply poset-independence of the underlying bundles
(Remark~\ref{rem:poset_indep}), which is what
Theorem~\ref{thm:aggregate_bound} Part~A requires.
\end{definition}

\subsection{Preservation of the lower bound}

\begin{proposition}[Commitment Preserves Part~A Lower Bound]
\label{prop:commit_preserves}
Theorem~\ref{thm:aggregate_bound} Part~A's aggregate lower bound
is computable from committed revealed-sacrifice events
(Definition~\ref{def:committed_event}) without access to the
hidden fields.  The bound depends only on the revealed fields
$(\Delta\volL(X_n), \catY[_n], t_n)$, all of which survive the
commitment's hiding property.
\end{proposition}

\begin{proof}[Proof sketch]
Part~A's lower bound is $\sum_{n=1}^{N} \Delta\volL(X_n)$ for
poset-independent bundles.  The sacrifice magnitudes
$\Delta\volL(X_n)$ are revealed directly.  Poset-independence is
verifiable from the revealed category labels: by
Definition~\ref{def:category_partition}, distinct labels
correspond to poset-independent partition components, and
containment is verified by
Definition~\ref{def:committed_event}(d).  No hidden field is
needed.
\end{proof}

\begin{remark}[Part~B requires additional committed fields]
Part~B's per-capability refinement requires $\alpha_{n,c}$
coefficients and stable capability identities across events.
The committed interface reveals only category labels.
For Part~B to operate on committed data, the commitment must
additionally expose a per-capability coordinate system within
each category component and the corresponding $\alpha$ values.
This is a richer commitment that reveals more structure than
Part~A requires.
\end{remark}

\begin{proposition}[ZK Aggregation with Poset-Independence]
\label{prop:zk_aggregation}
Zero-knowledge rollups over aggregate trade volume preserve the
lower-bound property \emph{including} the poset-independence
requirement.  Given a batch of $N$ committed sacrifice events,
a rollup proof can establish:
\begin{enumerate}
\item[(i)] $\sum_{n=1}^{N} \Delta\volL(X_n) \geq B$
  (aggregate bound);
\item[(ii)] the bundle category labels $\{\catY[_n]\}$ are
  pairwise distinct;
\end{enumerate}
for a public bound $B$, without revealing $N$, the individual
$\Delta\volL(X_n)$, or the specific category labels.

By Definition~\ref{def:category_partition}, distinct category
labels correspond to poset-independent partition components, so
(ii) establishes the poset-independence required by
Theorem~\ref{thm:aggregate_bound} Part~A.
\end{proposition}

\begin{proof}[Proof sketch]
The prover knows all $N$ committed events.  The rollup proof has
two components: (a)~a range proof on the aggregate sum
establishing $\sum \geq B$; (b)~a set-membership proof
establishing that the category labels are pairwise distinct,
using a Merkle commitment over the category set.
\end{proof}
